\clearpage
\section{사차방정식과 분해삼차식의 근호 구조}
\label{sec:quartic-radical-structure}

\begin{learninggoal}
$S_4$의 가해 정규열이 사차공식의 단계로 어떻게 나타나는지 설명한다. 우울형 사차식의 삼차분해식을 다시 구성하고, 그 세 근에서 제곱근을 취해 원래 네 근을 복원한다.
\end{learninggoal}

$\charac K\ne2,3$이라 하자. 일반 일계수 사차식
\[
  X^4+aX^3+bX^2+cX+d
\]
은
\[
  X=Y-\frac a4
\]
를 대입하여
\[
  f(Y)=Y^4+pY^2+qY+r
\]
꼴의 우울형 사차식으로 바뀐다.

일반적인 갈루아군 $S_4$에는 정규열
\[
  S_4\trianglerighteq A_4\trianglerighteq V_4
  \trianglerighteq C_2\trianglerighteq1
\]
이 있다. 인자군의 차수는 차례로
\[
  2,3,2,2
\]
이다. 따라서 체 쪽에서는 대략
\[
  \text{제곱근}\;\longrightarrow\;
  \text{세제곱근}\;\longrightarrow\;
  \text{제곱근}\;\longrightarrow\;
  \text{제곱근}
\]
의 구조가 예상된다. 실제 공식에서는 가운데 $S_4/V_4\cong S_3$ 작용을 하나의 삼차방정식으로 모으고, 그 삼차식을 Cardano 공식으로 푼 뒤 세 제곱근으로 네 근을 복원한다.

\subsection{세 쌍분할과 삼차분해식}

$f$의 네 근을 $x_1,x_2,x_3,x_4$라 하자. 우울형이므로
\[
  x_1+x_2+x_3+x_4=0.
\]
세 쌍분할에 대응하는 양을
\[
  z_1=(x_1+x_2)(x_3+x_4),
\]
\[
  z_2=(x_1+x_3)(x_2+x_4),
\]
\[
  z_3=(x_1+x_4)(x_2+x_3)
\]
로 둔다. $S_4$는 $z_1,z_2,z_3$을 순열하며, 이 작용의 핵은 $V_4$이다.

\begin{proposition}[우울형 사차식의 삼차분해식]
\label{prop:quartic-resolvent-cubic-radicals}
$z_1,z_2,z_3$은
\[
  \boxed{
  R_f(Z)=Z^3-2pZ^2+(p^2-4r)Z+q^2}
\]
의 세 근이다.
\end{proposition}

\begin{proof}
Vieta 관계는
\[
  \sum_i x_i=0,
  \qquad
  \sum_{i<j}x_ix_j=p,
  \qquad
  \sum_{i<j<k}x_ix_jx_k=-q,
  \qquad
  x_1x_2x_3x_4=r
\]
이다. 직접 정리하면
\[
  z_1+z_2+z_3=2p,
\]
\[
  z_1z_2+z_1z_3+z_2z_3=p^2-4r,
\]
\[
  z_1z_2z_3=-q^2.
\]
따라서 $z_i$들을 근으로 갖는 일계수 삼차식은 표시된 $R_f$이다.
\end{proof}

$R_f$의 분해체는 원래 분해체의 $V_4$-고정체이다. 따라서 삼차분해식을 푸는 것은
\[
  S_4\longrightarrow S_4/V_4\cong S_3
\]
인 몫작용을 해결하는 일이다. Cardano 공식 안의 제곱근과 세제곱근은 이 $S_3$ 층을 푼다.

\subsection{세 제곱근으로 네 근 복원하기}

$x_1+x_2+x_3+x_4=0$이므로
\[
  x_3+x_4=-(x_1+x_2)
\]
이고 따라서
\[
  z_1=-(x_1+x_2)^2.
\]
마찬가지로
\[
  z_2=-(x_1+x_3)^2,
  \qquad
  z_3=-(x_1+x_4)^2.
\]
제곱근의 부호를 적절히 택하여
\[
  a=\sqrt{-z_1}=x_1+x_2,
\]
\[
  b=\sqrt{-z_2}=x_1+x_3,
\]
\[
  c=\sqrt{-z_3}=x_1+x_4
\]
라 둘 수 있다. 그러면
\[
  a+b+c=2x_1
\]
이고 나머지 근도 부호를 바꾸어 얻는다.

\begin{proposition}[사차근의 복원]
\label{prop:quartic-root-recovery}
$z_1,z_2,z_3$이 삼차분해식 $R_f$의 세 근이라고 하자. 제곱근
\[
  a^2=-z_1,
  \qquad b^2=-z_2,
  \qquad c^2=-z_3
\]
를
\[
  abc=-q
\]
가 되도록 택하면 $f$의 네 근은
\[
  \boxed{x_1=\frac12(a+b+c)},
\]
\[
  \boxed{x_2=\frac12(a-b-c)},
\]
\[
  \boxed{x_3=\frac12(-a+b-c)},
\]
\[
  \boxed{x_4=\frac12(-a-b+c)}.
\]
\end{proposition}

\begin{proof}
위에서 $a=x_1+x_2$, $b=x_1+x_3$, $c=x_1+x_4$로 택하면 합과 차를 풀어 표시된 식을 얻는다. 부호조건은
\[
\begin{aligned}
abc
&=(x_1+x_2)(x_1+x_3)(x_1+x_4)\\
&=x_1^2(x_1+x_2+x_3+x_4)
 +\sum_{i<j<k}x_ix_jx_k\\
&=-q
\end{aligned}
\]
에서 나온다.

반대로 $a^2=-z_1$, $b^2=-z_2$, $c^2=-z_3$와 $abc=-q$를 만족하는 선택은 실제 근들에서 얻은 선택과 짝수개의 부호만 다르다. 표시된 네 조합은 순서만 달리하여 원래 네 근을 준다.
\end{proof}

\begin{theorem}[우울형 사차식의 근호해법]
\label{thm:quartic-radical-formula}
$\charac K\ne2,3$이고
\[
  f(X)=X^4+pX^2+qX+r
\]
라 하자. 먼저 삼차식
\[
  Z^3-2pZ^2+(p^2-4r)Z+q^2=0
\]
을 Cardano 공식으로 풀어 $z_1,z_2,z_3$을 구한다. 이어
\[
  a=\sqrt{-z_1},\qquad
  b=\sqrt{-z_2},\qquad
  c=\sqrt{-z_3}
\]
를 $abc=-q$가 되도록 택하면 명제 \ref{prop:quartic-root-recovery}의 네 식이 $f$의 모든 근을 준다.
\end{theorem}

\begin{pitfall}[제곱근의 부호]
$a,b,c$의 부호를 모두 독립적으로 택할 수는 없다. 조건
\[
  abc=-q
\]
가 한 개의 부호 자유도를 제거한다. 한 제곱근의 부호를 바꾸면 다른 하나의 부호도 함께 바꾸어야 같은 네 근의 순열을 얻는다.
\end{pitfall}

\subsection{판별식과 정규열의 해석}

삼차분해식의 판별식은 원래 사차식의 판별식과 같다.
\[
  \Disc(R_f)=\Disc(f).
\]
근 차이의 곱으로 보면
\[
\begin{aligned}
 z_1-z_2&=-(x_1-x_4)(x_2-x_3),\\
 z_1-z_3&=-(x_1-x_3)(x_2-x_4),\\
 z_2-z_3&=-(x_1-x_2)(x_3-x_4),
\end{aligned}
\]
이고 세 식의 제곱곱을 취하면 등식이 나온다.

일반적인 $S_4$의 경우 체의 층을 다음처럼 읽을 수 있다.
\[
\begin{array}{c|c|c}
\text{군의 단계} & \text{고정체의 단계} & \text{계산 도구}\\
\hline
S_4/A_4\cong C_2 & K\subset K(\sqrt\Delta) & \text{판별식의 제곱근}\\
A_4/V_4\cong C_3 & K(\sqrt\Delta)\subset L^{V_4} & \text{삼차분해식}\\
V_4/C_2\cong C_2 & L^{V_4}\subset L^{C_2} & \text{제곱근}\\
C_2/1\cong C_2 & L^{C_2}\subset L & \text{제곱근}
\end{array}
\]
실제 공식에서는 삼차분해식의 Cardano 풀이가 첫 두 층을 함께 처리하고, $\sqrt{-z_i}$들이 마지막 $V_4$ 층을 처리한다.

\begin{example}[짝다항식]
$q=0$이면
\[
  f(X)=X^4+pX^2+r
\]
이고 $Y=X^2$를 대입해 이차식
\[
  Y^2+pY+r=0
\]
으로 바로 줄일 수 있다. 삼차분해식도
\[
  Z\bigl(Z^2-2pZ+p^2-4r\bigr)
\]
로 인수분해된다. 이는 갈루아군이 일반적인 $S_4$보다 작아져 삼차층이 퇴화한 현상이다.
\end{example}

\begin{example}
\[
  X^4-2
\]
에서는 $p=q=0$, $r=-2$이므로 삼차분해식은
\[
  Z^3+8Z=Z(Z^2+8).
\]
분해체 $\Q(\sqrt[4]{2},i)$의 갈루아군은 $D_4$이며 가해군이다. 공식에 나타나는 제곱근들은 $\sqrt2$, $i$와 $\sqrt[4]{2}$를 차례로 복원한다.
\end{example}

\begin{keyidea}
사차공식이 삼차방정식을 먼저 푸는 이유는 $S_4$가 세 쌍분할에 작용하여
\[
  S_4/V_4\cong S_3
\]
를 만들기 때문이다. 삼차분해식은 이 몫작용을 기록하고, 나머지 $V_4$는 제곱근들로 해소된다.
\end{keyidea}

\begin{checkpoint}
\begin{enumerate}[label=(\alph*)]
  \item $X^4-5X^2+6$의 삼차분해식을 구하고 인수분해하라.
  \item $q=0$일 때 $abc=0$이 되는 이유를 삼차분해식의 상수항으로 설명하라.
  \item 사차식의 갈루아군이 $V_4$이면 삼차분해식이 기준체에서 완전히 분해되는 이유를 설명하라.
\end{enumerate}
\end{checkpoint}
